% 本文件由 LaTeX 排版 Agent 根据有效 translation.json 自动生成。
% author=Gennadius of Marseilles; work_id=2719; title=《名人传》补遗

\chapter{《名人传》补遗}

% structure: p0001 source=h1 target=omitted reason=page-title-already-rendered-by-parent

\Person{格纳迪乌斯}。\par

\Emph{在真福热罗尼莫去世后，格纳迪乌斯所增补的作者名单。}\par

% structure: p0004 source=h2 target=section reason=relative-heading-rank; base=h2

\section{第一章}

人称智者雅各伯者，乃波斯名城尼西比斯之主教，亦为在迫害者马克西米努斯手下之宣信者之一。他在尼西亚大公会议上，以反对之声推翻了亚略派关于\OriginalTerm{同质}{Homoousia}的乖谬。真福热罗尼莫虽在其《编年史》中提及此\Person{人}，赞其德行卓著，却未将其列入作者名录；此理易明，因我们注意到，他所提及的三四位叙利亚人，据他说是读过其希腊文译作的。由此可知，彼时他尚未通晓叙利亚语文或文献，故不识一位尚未被译成其他语言的作者。雅各伯的全部著作共有二十六卷，即：\WorkTitle{论信德}、\WorkTitle{驳一切异端}、\WorkTitle{论爱德}、\WorkTitle{论守斋}、\WorkTitle{论祈祷}、\WorkTitle{论对近人之特殊情感}、\WorkTitle{论复活}、\WorkTitle{论死后生命}、\WorkTitle{论谦卑}、\WorkTitle{论补赎}、\WorkTitle{论赔补}、\WorkTitle{论童贞}、\WorkTitle{论灵魂之价值}、\WorkTitle{论割损}、\WorkTitle{论受祝福的葡萄}、\WorkTitle{论依撒意亚之言：\BibleQuote{葡萄串必不毁坏}}、\WorkTitle{基督是天主子，与圣父同性同体}、\WorkTitle{论贞洁}、\WorkTitle{驳异民}、\WorkTitle{论会幕的建造}、\WorkTitle{论万民的归化}、\WorkTitle{论波斯王国}、\WorkTitle{论基督徒所受的迫害}。他还著有一部\WorkTitle{编年史}，对希腊人而言确实趣味不大，但可靠性极高，因其仅依据圣经之权威而构建。它堵住了那些凭狂妄臆测、闲散思辨关于敌基督或我主来临之人的口。雅各伯卒于君士坦提乌斯时代，按其父君士坦丁之命，葬于尼西比斯城墙之内，显然为保护该城；结果正如君士坦丁所料。多年后，尤利安进入尼西比斯，或嫉妒葬于此地者之荣耀，或嫉妒君士坦丁之信德——其家族因这种嫉妒而受迫害——遂下令将圣人之遗骸迁出城外。数月后，继尤利安之位的约维安皇帝，出于公共政策，将该城连同毗邻领土交给蛮族，该地至今受波斯统治。\par

% structure: p0006 source=h2 target=section reason=relative-heading-rank; base=h2

\section{第二章}

罗马主教尤利乌斯曾致信一位狄奥尼西，写有一封单篇书信\WorkTitle{论我主的降生成人}，在当时被认为有益于反驳那些宣称因降生而在基督内有两位格、因而亦有二性的人；但如今此信亦被视为有害，因为它助长了\OriginalTerm{欧提基派}{Eutychian}和\OriginalTerm{提摩太派}{Timothean}异端。\par

% structure: p0008 source=h2 target=section reason=relative-heading-rank; base=h2

\section{第三章}

保禄纳司铎，真福执事厄弗冷之弟子，为人精力充沛，精通圣经，在其师仍健在时，已为教会之杰出教师，尤其擅长即兴讲道。其师去世后，因贪图名誉，脱离教会，写下许多违背信仰之作。真福厄弗冷临终时，据说曾对站在身旁的保禄纳说：\Quote{你要当心，保禄纳，切勿顺从自己的思想；当你自以为完全认识天主时，要相信你尚未认识。}——因厄弗冷早已从保禄纳的钻研或言语中预感到，他正在探究新奇之事，并将其心智伸向无限；为此，厄弗冷常称其为新的巴尔代撒内。\par

% structure: p0010 source=h2 target=section reason=relative-heading-rank; base=h2

\section{第四章}

\Person{维特利乌斯}，非洲人，为多纳徒派\OriginalTerm{多纳徒派}{Donatist}分裂辩护，著有《论天主的仆人为何为世人所恨》，其中除了称我们为迫害者之外，发表了卓越的教义。他还著有\WorkTitle{驳异民}和驳斥我们这些在迫害时期背弃圣经者\OriginalTerm{背信者}{traditors}的作品，并写有大量\WorkTitle{论教会程序}。他在君士坦丁之子君士坦斯统治期间著名。\par

% structure: p0012 source=h2 target=section reason=relative-heading-rank; base=h2

\section{第五章}

\Person{玛克罗比乌斯}司铎，据我从奥普塔图斯的著作所知，后来在罗马秘密担任多纳徒派\OriginalTerm{多纳徒派}{Donatians}的主教。他写了一部作品\WorkTitle{致宣信者与童贞者}，此作在伦理上固然有益，且包含必要教义，并富有适用于保存贞洁的见解。他先在非洲属于我方，后来在罗马属于其本派，即多纳徒派或蒙丹派\OriginalTerm{蒙丹派}{Montanists}。\par

% structure: p0014 source=h2 target=section reason=relative-heading-rank; base=h2

\section{第六章}

\Person{赫略多鲁斯}司铎著有一书，题为\WorkTitle{论事物本质入门}，其中表明万物之源为一，无物与天主同永，天主并非恶的创造者，而是如此地创造一切善，以至于那用于恶的物质，也是在恶被发现后由天主所造；并且，任何物质性的东西都不能被认为是由天主之外的方式建立的，除天主之外别无创造者；天主以其预知知道本性将改变，遂以惩罚警告。\par

% structure: p0016 source=h2 target=section reason=relative-heading-rank; base=h2

\section{第七章}

\Person{帕科米乌}隐修士，一位在教导和行奇迹上都拥有宗徒恩宠的人，也是埃及隐修院的创建者，撰写了一部\WorkTitle{纪律秩序}，适合两类隐修士，是他由天使口授而得的。他也写信给他教区内的联合主教们，用一种被神秘圣事所隐藏的字母书写，以超越惯常的人类知识，唯有那些拥有特殊恩宠或功德者才能明瞭，即：一封\WorkTitle{致科尔内略院长}、一封\WorkTitle{致叙禄院长}，以及一封\WorkTitle{致所有隐修院院长}，劝勉他们聚集到一座非常古老的隐修院——其埃及语称为\Quote{\Place{鲍}}——共同以永久的法律庆祝逾越节。他在另一封信中也敦促，在八月庆祝的赎罪之日，各首席主教应聚集在一处。他还写了另一封信给被派到隐修院外工作的弟兄们。\par

% structure: p0018 source=h2 target=section reason=relative-heading-rank; base=h2

\section{第八章}

\Person{德奥多鲁斯}，上述院长帕科米乌在恩宠与领导权上的继任者，致信其他隐修院，信件以圣经的语言写成，然而他在信中频繁提及他的导师与老师帕科米乌，并以他的教义与生活为榜样。他说，这些是他由一位天使所教导，好让自己也能再教导别人。他同样劝勉他们持守内心与意愿的目的，并使那些在院长去世后发生纷争、因脱离团体而破坏了合一的人重归于和谐与合一。他的三封劝勉书信现存。\par

% structure: p0020 source=h2 target=section reason=relative-heading-rank; base=h2

\section{第九章}

\Person{奥雷西西斯}隐修士，帕科米乌与德奥多鲁斯两人的同工，一位在圣经上学至完美的人，撰写了一部书，以神盐调味，并由所有隐修纪律的精要构成——中庸地说——其中几乎整部旧约与新约都以紧凑的论述呈现出来，至少是关乎隐修士特别需要的一切。他将此书交给他的弟兄们，几乎就在他去世的那一天，如同留下了一份遗产。\par

% structure: p0022 source=h2 target=section reason=relative-heading-rank; base=h2

\section{第十章}

\Person{玛加略}，埃及隐修士，以奇迹与德行著称，写了一封信，是致同修中的较年轻者。在此信中，他教导他们说：唯有那认识自身受造状况，并献身于一切劳苦，且与今生一切悦人事物搏斗，同时恳求天主助佑的人，才能完美地侍奉天主，并获致本性纯洁、获得节制，作为理应得到的本性恩赐。\par

% structure: p0024 source=h2 target=section reason=relative-heading-rank; base=h2

\section{第十一章}

\Person{埃瓦格里乌}隐修士，上述玛加略的亲密门徒，受过圣学与俗学教育，为人卓异。那部名为\WorkTitle{教父生平}的书提及他是一位极其节制和博学的人。他撰写了许多对隐修士有益的作品，其中包括：\WorkTitle{驳八种主要邪念建议}。他是最早，或至少是首批教导这些邪念的人，并针对它们撰写了八部书，仅从圣经的见证中取材，效法吾主的榜样，祂总是以圣经的引文来应对试探者，以便每一邪念，无论来自魔鬼还是来自败坏的本性，都有一见证加以驳斥。这部作品我受指示翻译成了拉丁文，以我在希腊文中发现的同样朴素风格翻译。他还撰写了一部\WorkTitle{百感}，供那些单纯独修者使用，分章编排；以及一部\WorkTitle{五十感}，供给博学和勤学者，这是我首次翻译成拉丁文。前者先前已被翻译，我通过重译和订正部分内容将其恢复，以呈现作者的真实意义，因为我看到译作因时间而遭败坏和混淆。他还撰写了一部适合于会院修士和团体修士及奉献于天主的贞女的公共生活教义，是一本适合于其宗教与性别的小书。他也出版了几本极为隐晦的意见集，据他自己说，只有隐修士的心才能领悟，这些我也以拉丁文出版。他寿至老年，大显征兆与奇迹。\par

% structure: p0026 source=h2 target=section reason=relative-heading-rank; base=h2

\section{第十二章}

\Person{德奥多鲁斯}，安提约基雅教会的司铎，一位审慎的探究者且口才伶俐，著书反驳\OriginalTerm{阿波里纳异端者}{Apollinarians}和\OriginalTerm{阿诺米异端者}{Anomians}，书名\WorkTitle{论主的降生成人}，共十五卷，包含一万五千句诗，他以最清晰的推理和圣经的见证证明：主耶稣如何具有圆满的天主性，同样也具有圆满的人性。他还教导说，人仅由两种实体组成：灵魂和肉身；感觉和心神并非不同的实体，而是灵魂内在固有的官能，灵魂借此受到灵感并具有理性，并借此使肉身拥有感觉能力。此外，这部作品的第十四卷完全论述天主圣三的非受造、唯一无形且宰制的本性，以及动物的理性，他以虔诚的精神、凭借圣经的权威加以解释。在第十五卷，他通过引用教父的传统来确认和巩固他整部作品的整体。\par

% structure: p0028 source=h2 target=section reason=relative-heading-rank; base=h2

\section{第十三章}

\Person{普鲁登修斯}，一位精通世俗文学的人，撰写了一部\WorkTitle{特罗凯乌姆}，选自整部旧约与新约的精选人物。他还按照希腊人的方式撰写了一部注释\WorkTitle{论创世六日}，从世界的创造直到第一个人的创造及其堕落。他还撰写了几部短书，希腊文标题为\WorkTitle{神化}、\WorkTitle{灵性之战}和\WorkTitle{罪之起源}，即\WorkTitle{论天主性}、\WorkTitle{论灵性冲突}、\WorkTitle{论罪的起源}。他还撰写了\WorkTitle{殉道者颂}，一部书邀请人殉道，引用数位殉道者为例；以及另一部\WorkTitle{圣歌集}，但特别针对为偶像崇拜辩护的\WorkTitle{驳叙玛古斯}，从该书中我们得知\Person{帕拉丁}是一名士兵。\par

% structure: p0030 source=h2 target=section reason=relative-heading-rank; base=h2

\section{第十四章}

\Person{Audentius}，西班牙主教，写了一本书反对摩尼教徒、撒伯流派和亚流派，更特别反对福提诺派，即现在称为波诺西派的。他将这本书命名为\WorkTitle{论信德反对异端}，并在其中证明圣子与圣父同永，并且圣子并非在因天主的行动而受孕、由母亲童贞玛利亚以真实人性出生时，才从天主圣父获得其神性的开始。\par

% structure: p0032 source=h2 target=section reason=relative-heading-rank; base=h2

\section{第十五章}

\Person{Commodianus}，在从事世俗文学时也阅读我们的著作，并找到机会接受了信仰。如此成为基督徒后，他愿将自己的研究成果献给基督——他救恩的作者，便以勉强可忍受的半韵文语言写了一部\WorkTitle{反对异教徒}的书，并且由于他对我们的文学知之甚少，他更能推翻他们的教义，而非建立我们的。因此，在关于神圣应许的争论中，他以一种相当可怜且可以说是粗糙的方式谈论，令他们惊愕，令我们绝望。他以\Person{特土良}、\Person{拉克坦提乌斯}和\Person{帕皮亚}为权威，采纳并教导他的学生良好的伦理原则，特别是自愿爱贫穷。\par

% structure: p0034 source=h2 target=section reason=relative-heading-rank; base=h2

\section{第十六章}

司铎\Person{Faustinus}写信给皇后\Person{Flaccilla}，写了七卷书\WorkTitle{反对亚流派和马其顿派}，引用他们所用以亵渎的圣经本身的见证，进行论证并定罪他们。他还与一位名叫\Person{Marcellinus}的司铎共同写了一本书，致皇帝\Person{Valentinianus}、\Person{Theodosius}和\Person{Arcadius}，为他们的同道基督徒辩护。从这本书可以看出他顺从了路济弗尔裂教，因为在这本书中他责备普瓦捷的\Person{希拉利}和罗马主教\Person{Damasus}，指责他们给了教会不当的忠告，建议为了恢复和平而接纳背教的主教。因为路济弗尔派不愿接纳那些在阿里米诺会议中曾与亚流共融的主教，正如诺洼天派不愿接纳悔改的背教者一样。\par

% structure: p0036 source=h2 target=section reason=relative-heading-rank; base=h2

\section{第十七章}

\Person{鲁菲努斯}，阿奎莱亚教会司铎，在教会的圣师中并非最小，且具有将希腊文优雅译成拉丁文的卓越才能。以此方式，他为拉丁教会打开了希腊文学的大部分：翻译了凯撒勒亚的\Person{巴西略}、最善辞令的纳齐盎的\Person{格里高利}、罗马的\Person{克莱孟}的\WorkTitle{认识篇}、巴勒斯坦凯撒勒亚的\Person{优西比乌}的\WorkTitle{教会史}、\Person{西克斯图斯}的\WorkTitle{语录}、\Person{埃瓦格里乌斯}的\WorkTitle{语录}以及殉道者\Person{庞斐禄}的著作\WorkTitle{反对数学家}。在这些被拉丁读者阅读的作品中，凡有序言的，都是鲁菲努斯所译；而没有序言的，则是别人所译，那人未曾选择写序言。然而，并非所有奥利振的作品都是他所译，因为\Person{热罗尼莫}翻译了一些，由他的序言可辨。同样的鲁菲努斯，藉天主的恩宠，自己出版了一部如此卓越的\WorkTitle{宗徒信经释义}，以至于其他释义与之相比都算不得什么。他还以三重意义（即历史、道德和神秘意义）论述了雅各对族长的祝福。他还写了许多劝勉敬畏天主的书信，其中写给\Person{Proba}的尤为出众。他还为上述由优西比乌所写且由他翻译的教会史增加了第十卷和第十一卷。此外，他回应了对他作品的诽谤者，用了两卷书的篇幅，论证并证明他是在主的助佑下、在天主面前，为教会的益处而运用才能，而对方则是受嫉妒所激，投入了争论。\par

% structure: p0038 source=h2 target=section reason=relative-heading-rank; base=h2

\section{第十八章}

\Person{Tichonius}，非洲人，据称对圣经文学足够精通，对世俗文学也并非完全陌生，且热衷教会事务。他写了\WorkTitle{论内部战争}和\WorkTitle{各种原因的阐发}，在其中为辩护他的朋友而引用古代会议，从这一切可看出他是个多纳徒派。他还撰写了八条\WorkTitle{探求并确定圣经意义的规则}，压缩成一卷。他还完整阐释了若望默示录，认为其中没有丝毫属肉体的意义，全是属灵的。在这阐释中，他主张天使的本性是有形体的；此外，他怀疑义人复活后在地上统治一千年，或者有两次肉身复活——一次义人、一次不义者；但他主张所有将有一次同时的复活，届时连流产的胎儿和畸形儿也要复活，免得任何活着的人，无论多么畸形，会失落。他确实在两种复活之间作了区分，将第一种称为义人的启示，仅发生在教会的成长中，在那里因信德成义的人通过圣洗从罪恶的死亡身体中复活，服务于永生；而第二种是所有人肉身的一般复活。此人活跃于上述\Person{鲁菲努斯}同时期，在\Person{德奥多西}皇帝及其儿子执政期间。\par

% structure: p0040 source=h2 target=section reason=relative-heading-rank; base=h2

\section{第十九章}

\Person{塞维鲁}司铎，别号\Person{苏尔皮基乌斯}，出身\Place{阿基坦}行省，因出身高贵、文学成就卓越、甘于神贫与谦卑而闻名，亦为圣德之人图尔主教圣\Person{玛尔定}与诺拉的圣\Person{保利诺}所爱慕。他写了一些毫不低劣的小书。他写给其妹妹许多\WorkTitle{书信}，劝勉人爱慕天主、轻视世俗；这些书信广为流传。他写给上述诺拉的\Person{保利诺}两封信，也写给其他人一些信，但因其中几封涉及家事，故未收集出版。他还撰写了一部\WorkTitle{编年史}，也为许多人益处写了\WorkTitle{圣玛尔定传}，这位隐修士及主教是位以奇事、奇迹与德行闻名之人。他还写作了\WorkTitle{波斯图米亚努斯与加卢斯对话录}，他在其中充当辩论的调解人与裁判。主题涉及东方隐修士及圣玛尔定的生活方式——一种分为两部分的对话。在第一部分中，他提到在他那个时代\Place{亚历山大里亚}教务会议上的一项主教法令，大意是：有智慧者应阅读\Person{奥力振}，虽须谨慎，因其内有善；但能力较弱者应因其恶而弃绝之。晚年，他被白拉奇主义者引入歧途，又认识到多言有罪，遂保持缄默直到逝世，以期通过忏悔的沉默赎清因言语所犯之罪。\par

% structure: p0042 source=h2 target=section reason=relative-heading-rank; base=h2

\section{第二十章}

\Person{安提约基}主教写了一部长篇著作\WorkTitle{驳贪婪}，还编写了一篇充满虔敬悔改与谦卑的讲道词\WorkTitle{论救主医治盲人复明}。他在皇帝\Person{阿卡迪乌斯}在位期间去世。\par

% structure: p0044 source=h2 target=section reason=relative-heading-rank; base=h2

\section{第二十一章}

\Person{塞维里亚努斯}，\Place{加巴拉}教会主教，精通圣经，且是出色的宣道者。为此，主教\Person{若望}和皇帝\Person{阿卡迪乌斯}经常召他在\Place{君士坦丁堡}讲道。我读过他的\WorkTitle{迦拉达书释义}和一部非常吸引人的小作品\WorkTitle{论圣洗与主显节}。他在其洗礼之子\Person{德奥多西}在位期间去世。\par

% structure: p0046 source=h2 target=section reason=relative-heading-rank; base=h2

\section{第二十二章}

\Person{尼凯亚斯}，\Place{罗马提亚}城主教，以简洁明了的语言撰写了六卷\WorkTitle{新领洗者教导}。第一卷论如何寻求获得圣洗恩宠的候洗者应当如何行事；第二卷论亲属关系之错误，其中提到，在他自己时代不久前，某一家族之父\Person{梅洛迪乌斯}因慷慨，农民\Person{加拉迪乌斯}因勇敢，被异教徒列于诸神之中；第三卷\WorkTitle{论独一主宰的信德}；第四卷\WorkTitle{驳祖谱}；第五卷\WorkTitle{论信经}；第六卷\WorkTitle{论逾越节羔羊祭献}。他还撰写了\WorkTitle{给堕落的贞女}，激励所有堕落者悔改。\par

% structure: p0048 source=h2 target=section reason=relative-heading-rank; base=h2

\section{第二十三章}

\Person{奥林皮乌斯}主教，西班牙裔，撰写了一卷论信德的书，驳斥那些指责本性而非意志之人，指出恶并非由创造引入本性，而是由悖逆引入。\par

% structure: p0050 source=h2 target=section reason=relative-heading-rank; base=h2

\section{第二十四章}

\Person{巴基亚里乌斯}，一位基督宗教哲学家，敏捷而果断，决心将时间奉献给天主，选择旅行作为保全其意向纯全的方法。据说他出版了一些受人欢迎的小作品，但我只读过其中一部，即\WorkTitle{论信德}，他在其中向该城首席司铎为自己陈辩，对抗那些抱怨并误传其旅行的人，声称他旅行并非出于对人的恐惧，而是为了天主的缘故，从他本族和亲族中出来，成为宗祖\Person{亚巴郎}的共同继承人。\par

% structure: p0052 source=h2 target=section reason=relative-heading-rank; base=h2

\section{第二十五章}

\Person{萨巴提乌斯}，\Place{高卢}行省主教，应一位名为\Person{塞昆达}的贞女（她贞洁且奉献于基督）之请，撰写了一卷\WorkTitle{论信德}反驳\Person{玛尔基翁}及其师傅\Person{瓦伦提努斯}，也反驳\Person{尤诺米乌斯}及其师傅\Person{埃提乌斯}，通过理性与圣经的证词显示：天主性的来源是唯一的，祂的永恒的作者与从无中创造大地的创造者是同一位；关于基督也是如此：祂并非以幻象显现为人，而是具有真实血肉，通过这血肉——饮食、疲倦、哭泣、受苦、死亡、复活——显示祂确实是人。因为\Person{玛尔基翁}和\Person{瓦伦提努斯}反对这些意见，主张天主性来源是双重的，基督以幻象出现。对于\Person{埃提乌斯}及其弟子\Person{尤诺米乌斯}，他指明：圣父与圣子并非具有两种天性且在天主性上平等，而是同一实体，一位出自另一位，即圣子出自圣父，相互永存；这是埃提乌斯与尤诺米乌斯所反对的信理。\par

% structure: p0054 source=h2 target=section reason=relative-heading-rank; base=h2

\section{第二十六章}

\Person{依撒格}撰写了\WorkTitle{论天主圣三}和\WorkTitle{论主的降生成人}，以非常晦涩的论证风格和纠缠的语言写作，主张在一位天主性中存在三个位格，以至于每个位格拥有其他位格所不具有的特性：即圣父的特性是非由谁而生，是祂人的起源；圣子的特性是受生而非在生者之后；圣神的特性是非受造亦非受生，却出自另一位。关于主的降生成人，他写道：天主圣子的位格被认为是一个，而其中却有两个本性存在。\par

% structure: p0056 source=h2 target=section reason=relative-heading-rank; base=h2

\section{第二十七章}

\Person{乌尔西努斯}隐修士撰文反驳那些主张异端者应重洗之人，教导说这是不合法且不尊崇天主的——那些曾受洗者，无论是只奉基督之名受洗，或奉圣父、圣子、圣神之名受洗，即使所用公式意义不正，也不应重洗。他认为，在简单的宣认天主圣三与基督之后，天主教司铎的覆手便足以获得救恩。\par

% structure: p0058 source=h2 target=section reason=relative-heading-rank; base=h2

\section{第二十八章}

\Person{玛加略} 另一位隐修士，在罗马写了\WorkTitle{反对数学家}一书，在这项工作中他寻求东方著作的慰藉。\par

% structure: p0060 source=h2 target=section reason=relative-heading-rank; base=h2

\section{第二十九章}

\Person{赫略多鲁斯}，安提约基雅司铎，发表了一部从圣经中汇集的优秀著作\WorkTitle{论童贞}。\par

% structure: p0062 source=h2 target=section reason=relative-heading-rank; base=h2

\section{第三十章}

[ \Person{若望}，君士坦丁堡主教，一个拥有奇妙知识和圣洁生活的人，在各样事上都值得效法，为所有奔向神圣事物的人写了许多非常有用的著作。其中有\WorkTitle{论灵魂的痛悔}一书、\WorkTitle{除了自己无人能伤害他}、一部优秀著作\WorkTitle{赞美有福的宗徒保禄，论近卫总管\Person{欧特罗比乌斯}的过分行为和恶名}以及许多其他著作，如我所说，勤奋者可以找到它们。]\par

% structure: p0064 source=h2 target=section reason=relative-heading-rank; base=h2

\section{第三十一章}

\Person{另一位若望}，耶路撒冷主教，写了一本书反对那些贬低他学问的人，在书中他表明他追随\Person{奥力振}的才智而非他的信经。\par

% structure: p0066 source=h2 target=section reason=relative-heading-rank; base=h2

\section{第三十二章}

\Person{保禄}主教写了一部短著\WorkTitle{论补赎}，他在书中为补赎者制定这条律则：他们应当以这样的方式为自己的罪悔改，以至于不会被绝望的悲伤过度压倒。\par

% structure: p0068 source=h2 target=section reason=relative-heading-rank; base=h2

\section{第三十三章}

\Person{赫耳维狄乌}，是\Person{奥克森狄}的门徒、\Person{西玛库斯}的模仿者，他写作时确实怀着宗教热忱，但不按知识，写了一本语言和推理均不精炼的书。在这部著作中，他如此企图把圣经的意思扭曲成自己的乖谬，以致敢于依据它们断言：\Person{若瑟}和\Person{玛利亚}在我们主诞生之后有了孩子，被称为主的兄弟。为回应他的乖谬，热罗尼莫发表了一本反对他的书，充满圣经证据。\par

% structure: p0070 source=h2 target=section reason=relative-heading-rank; base=h2

\section{第三十四章}

\Person{德奥菲禄}，亚历山大里亚教会的主教，写了一卷巨著\WorkTitle{反对奥力振}，在其中他几乎谴责了奥力振的所有言论和他本人，同时说他（奥力振）的见解并非独创，而是从古代教父、尤其从\Person{赫拉克拉斯}那里得来，说他被免去司铎职务，被逐出教会，被迫逃离城市。他还写了\WorkTitle{反对拟人论者}，拟人论者是那些说天主有人形和肢体的异端，他用长篇讨论和圣经的见证进行驳斥和论证，令人信服。他显示，按照教父们的信仰，天主应被想作无形的，绝无任何肢体的暗示，因此在受造物中没有任何东西在本质上像祂；祂本性的不朽、不变、无形也没有赐给任何人，但所有理智本性都是有形的、可朽的、可变的，只有祂不应受朽坏或改变，只有祂拥有不朽和生命。同样，尼西亚大公会议发现逾越节庆典在九十年后会在同一时间、同一月、同一天复现，他加上了关于节庆的一些观察和解释，呈给\Person{戴奥多西}皇帝。我也读过三卷\WorkTitle{论信德}的书，署名是他，但文风不像他的，我不太认为是他的作品。\par

% structure: p0072 source=h2 target=section reason=relative-heading-rank; base=h2

\section{第三十五章}

\Person{欧瑟比}写了\WorkTitle{论我们主十字架的奥迹}和宗徒们的忠信，尤其\Person{伯多禄}由十字架德能所获得的。\par

% structure: p0074 source=h2 target=section reason=relative-heading-rank; base=h2

\section{第三十六章}

\Person{维吉兰狄}，高卢公民，掌管巴塞罗那教会。他也写作，带着一些宗教热忱，但被追求人称赞的欲望所胜，且自不量力；他是一个语言精炼但未在圣经意义上操练的人。他以歪曲的意义解释\Person{达尼尔}的神视，并说了其他轻浮的事，这些在异端名录中必须提及。[有福的司铎热罗尼莫也回应了他。]\par

% structure: p0076 source=h2 target=section reason=relative-heading-rank; base=h2

\section{第三十七章}

\Person{辛普利齐亚诺}，主教，在多封信中劝勉当时为司铎的奥斯定，他应当运用自己的才智，花时间解释圣经，如此，一个类似新\Person{盎博罗削}、\Person{奥力振}的监督者可能出现。因此他也送给他许多圣经考题。还有一封他的书信\WorkTitle{问题集}，在其中他以问问题的方式教导，好像希望学习一样。\par

% structure: p0078 source=h2 target=section reason=relative-heading-rank; base=h2

\section{第三十八章}

\Person{维吉禄}主教写给一位\Person{辛普利齐亚诺}一本小书\WorkTitle{赞美殉道者}以及一封包含他时期在蛮族中殉道者行传的书信。\par

% structure: p0080 source=h2 target=section reason=relative-heading-rank; base=h2

\section{第三十九章}

\Person{奥斯定}，非洲人，希波勒吉斯主教，因神圣和世俗的学识而闻名于世，信德无瑕，生活纯洁，著作如此之多，以致不能全部收集。因为谁能自夸拥有他全部著作，或谁如此勤奋读完他写的一切？甚至在老年时，他发表了十五卷\WorkTitle{论天主圣三}，这是他在年轻时就开始写的。在其中，正如圣经所说，被带入君王的宝库，以天主智慧多彩的衣裳为装饰，他展示了一个没有瑕疵、没有皱纹或任何此类东西的教会。在他的著作\WorkTitle{论主的降生成人}中，他也表现了一种独特的虔诚。关于死者复活，他以同样的真诚写作，而将关于流产的疑问留给能力较弱的人。\par

% structure: p0082 source=h2 target=section reason=relative-heading-rank; base=h2

\section{第四十章}

\Person{奥罗修}，一位西班牙司铎，一个极其雄辩、精通历史的人，写了八卷书反对那些基督徒的敌人，他们声称罗马帝国的衰落是由基督徒的宗教造成的。在这些书中，他叙述了从创世以来几乎整个世界的灾难、痛苦和战争骚乱，表明罗马帝国应归功于基督徒的宗教而获得了本不应得的延续，以及它为敬拜天主所享有的和平状态。\par

在第一卷中，他描述了位于永流大洋（俄刻阿诺斯）之中、被塔奈斯河（顿河）分割的世界，记载了各地位置、各民族名称、数目及风俗、各地区的特征、引发的战争以及以亲友鲜血凝固而成的帝国形成。\par

这就是那位奥罗修斯，他受奥斯定派遣，往热罗尼莫处学习灵魂之性；返回时，最先给西方带来了新近发现的真福首殉者斯德望的圣髑。他活跃于霍诺里乌斯皇帝统治末期。\par

% structure: p0086 source=h2 target=section reason=relative-heading-rank; base=h2

\section{第四十一章}

\Person{玛克西穆斯}，都灵教会主教，研读圣经颇为勤勉，且擅长即席教导民众。他撰写了\WorkTitle{宗徒与洗者若翰颂}、\WorkTitle{诸殉道者讲道}。此外，他还对福音和宗徒大事录的多处经文作了精辟注释。他还撰写了\WorkTitle{圣欧瑟伯传}（维切利主教、宣信者）、\WorkTitle{圣西彼廉论}，并发表了\WorkTitle{圣洗恩宠论}单行本。我读过他的以下作品：\WorkTitle{贪吝论}、\WorkTitle{款待论}、\WorkTitle{月食论}、\WorkTitle{施舍论}、\WorkTitle{依撒意亚预言：你的酒商掺水论}、\WorkTitle{吾主受难论}、\WorkTitle{天主仆役守斋通论}、\WorkTitle{四旬斋论}、\WorkTitle{守斋日不可戏谑论}、\WorkTitle{犹达斯论}（背叛者）、\WorkTitle{吾主十字架论}、\WorkTitle{吾主坟墓论}、\WorkTitle{吾主复活论}、\WorkTitle{吾主在般雀比拉多前受控审讯论}、\WorkTitle{一月朔日论}、\WorkTitle{吾主圣诞讲道}、\WorkTitle{主显节讲道}、\WorkTitle{逾越节讲道}、\WorkTitle{圣神降临讲道}，还有许多\WorkTitle{勿惧血肉之敌论}、\WorkTitle{饭后谢恩论}、\WorkTitle{尼尼微人悔改论}，以及其他他在不同场合发表的讲道，其名称我已不记得。他去世于霍诺里乌斯与幼帝狄奥多西统治时期。\par

% structure: p0088 source=h2 target=section reason=relative-heading-rank; base=h2

\section{第四十二章}

\Person{伯多尼}，意大利波洛尼亚主教，圣德昭著，自幼修习隐修之道。据传他撰写了\WorkTitle{教父列传}，即埃及隐修士列传，隐修士们视之为其修道的明镜与模范。我读过一部题为\WorkTitle{主教祝圣论}的论著，理据充分，谦逊卓著，但其文笔典雅，表明非其所作；或如有些人所说，为其父伯多尼之作，其人极擅辞令，饱学世俗文学。我以为此说可接受，因该著作作者自称为近卫总领。他去世于狄奥多西与瓦伦提尼安统治时期。\par

% structure: p0090 source=h2 target=section reason=relative-heading-rank; base=h2

\section{第四十三章}

\Person{佩拉纠}，异端魁首，在被宣布为异端者之前，撰写了若干对学生有实用价值的作品：\WorkTitle{信仰天主圣三论}三卷，及\WorkTitle{圣经论基督徒生活选录}一卷。后者附有目录，仿效殉道者圣西彼廉的体例。然而，在被宣布为异端者之后，他撰写了与其异端相关的著作。\par

% structure: p0092 source=h2 target=section reason=relative-heading-rank; base=h2

\section{第四十四章}

\Person{依诺增爵}，罗马主教，撰写了西方教会针对佩拉纠派所通过的法令；其继承者教宗佐西马随后广为颁行。\par

% structure: p0094 source=h2 target=section reason=relative-heading-rank; base=h2

\section{第四十五章}

\Person{塞莱斯提乌}，在加入佩拉纠之前，尚属年轻时，致书其父母，写了三封\WorkTitle{论隐修生活}的书简，皆以短书形式写成，包含适合每位寻求天主之人的道德箴言，毫无后来显现的过失痕迹，全篇致力于劝勉修德。\par

% structure: p0096 source=h2 target=section reason=relative-heading-rank; base=h2

\section{第四十六章}

\Person{儒利安}主教，性格刚毅，精通圣经，兼擅希腊文与拉丁文；在暴露其参与佩拉纠之邪说以前，在教会圣师中卓然超群。但此后，他力图维护佩拉纠异端，撰写了\WorkTitle{驳奥斯定}四卷（奥斯定乃佩拉纠之反对者），随后又有八卷。另有一卷讨论集，双方各陈己见。\par

这位儒利安，在饥馑匮乏之时，以其布施之慷慨及所施美德之魅力，吸引了许多人，使他们与他一同陷入异端。他去世于君士坦提乌斯之子瓦伦提尼安统治时期。\par

% structure: p0099 source=h2 target=section reason=relative-heading-rank; base=h2

\section{第四十七章}

\Person{路济安}司铎，一位圣德之人。在霍诺里乌斯与狄奥多西为帝时，天主启示了他圣斯德望\OriginalTerm{首殉者}{Protomartyr}坟墓及遗骸之所在。他以希腊文记录了那启示，致信所有教会。\par

% structure: p0101 source=h2 target=section reason=relative-heading-rank; base=h2

\section{第四十八章}

\Person{阿维图}司铎，西班牙裔，将上述路济安司铎的著作译为拉丁文，并附上自己的信函，委托司铎奥罗修斯转交西方教会。\par

% structure: p0103 source=h2 target=section reason=relative-heading-rank; base=h2

\section{第四十九章}

\Person{保利诺}，坎帕尼亚诺拉主教，创作了许多简短的韵文作品，以及一部致切尔苏斯的\WorkTitle{慰藉书\newline{}论基督徒受洗孩童之死}，近乎墓志铭，充满基督徒的望德；还有许多\WorkTitle{致塞维鲁书简}，以及他在就任主教前所写的散文\WorkTitle{颂词：战胜暴君论}，呈送给狄奥多西，主张胜利在于信德与祈祷，而非武器。他还编纂了\WorkTitle{圣事礼书}和\WorkTitle{圣歌集}。\par

他还致函其姐妹，论及\WorkTitle{蔑世论}，并在不同场合发表了各种类型的论著。\par

他所有小作品中最著名的，当数\WorkTitle{悔改论}和\WorkTitle{诸殉道者总颂}。他生活在霍诺里乌斯与瓦伦提尼安统治时期，不仅以其博学与圣德著称，还因其驱魔能力而闻名。\par

% structure: p0107 source=h2 target=section reason=relative-heading-rank; base=h2

\section{第五十章}

司铎\Person{欧特罗皮乌斯}致信两位姊妹，即基督的婢女，她们因对贞洁的热忱和对宗教的爱好而被父母剥夺继承权。他以小册子的形式写了《慰藉书信》两卷，行文优美清晰，不仅以论证，也以圣经的见证加以巩固。\par

% structure: p0109 source=h2 target=section reason=relative-heading-rank; base=h2

\section{第五十一章}

另一位\Person{埃瓦格里乌斯}写了一部\WorkTitle{犹太人西蒙与基督徒德奥斐卢斯之间的讨论}，这部作品非常著名。\par

% structure: p0111 source=h2 target=section reason=relative-heading-rank; base=h2

\section{第五十二章}

执事\Person{维基利乌斯}根据教父的传统编撰了\WorkTitle{隐修士规则}，该规则常在隐修院中被诵读，以裨益聚集的隐修士。它以简洁清晰的语言写成，涵盖了隐修生活的全部职责。\par

% structure: p0113 source=h2 target=section reason=relative-heading-rank; base=h2

\section{第五十三章}

君士坦丁堡主教\Person{阿提库斯}致信阿尔卡狄乌斯皇帝的公主女儿们，写了\WorkTitle{论信德与童贞}，一部非常卓越的作品，他在其中预先攻击了聂斯多略的教理。\par

% structure: p0115 source=h2 target=section reason=relative-heading-rank; base=h2

\section{第五十四章}

异端之首\Person{聂斯多略}，当他在安提约基雅教会担任司铎时，被视为一位杰出的即兴教师，并撰写了大量关于各种\WorkTitle{问题}的论文，当时他已经将那种微妙的恶注入其中，后来这种恶成了公认不虔敬的毒素，同时却以道德劝诫为遮掩。但后来，因其口才和节制而被推举为君士坦丁堡教会的主教，他公开显露出长久以来隐藏的东西，成为教会的公开敌人，并写了一本书\WorkTitle{论主的降生成人}，由六十二段神圣圣经的经文组成，这些经文被曲解使用。他在此书中所坚持的内容，可在异端者目录中找到。\par

% structure: p0117 source=h2 target=section reason=relative-heading-rank; base=h2

\section{第五十五章}

罗马主教\Person{塞勒斯提努斯}致函东西方教会，报告了公会议反对上述聂斯多略的法令，并主张在基督内虽有两个完整的性体，但天主圣子的位格应视为单一。上述聂斯多略被证明反对这一观点。同样，塞勒斯提努斯的继任者\Person{西斯图斯}也就同一主题致信同一聂斯多略和东方主教们，陈述了西方主教们反对其错误的观点。\par

% structure: p0119 source=h2 target=section reason=relative-heading-rank; base=h2

\section{第五十六章}

加拉提亚地区安基拉的主教\Person{德奥多图斯}，在厄弗所时，撰写了反对聂斯多略的辩护与驳斥之作，诚然以辩证风格写成，但交织着圣经的经文。他的方法是先作出陈述，然后从圣经中引用佐证经文。\par

% structure: p0121 source=h2 target=section reason=relative-heading-rank; base=h2

\section{第五十七章}

不列颠的主教\Person{法斯蒂迪乌斯}致信一位名叫法塔利斯的人，写了一本\WorkTitle{论基督徒生活}，另一本\WorkTitle{论保守童贞的境况}，这是一部充满正确教理并光荣天主的作品。\par

% structure: p0123 source=h2 target=section reason=relative-heading-rank; base=h2

\section{第五十八章}

亚历山大里亚教会的主教\Person{济利禄}发表了关于各种\WorkTitle{问题}的多篇论文，并撰写了许多讲道集，这些讲道被希腊主教们推荐用于讲道。他的其他著作有：\WorkTitle{论会堂的倾覆}、\WorkTitle{论信德——驳斥异端者}，以及一部特别针对聂斯多略的作品，标题为\WorkTitle{驳斥}，其中揭露了聂斯多略的一切秘密，并驳斥了他已发表的观点。\par

% structure: p0125 source=h2 target=section reason=relative-heading-rank; base=h2

\section{第五十九章}

主教\Person{提摩太}写了一本书\WorkTitle{论吾主按照肉身的诞生}，据信是在主显节写成的。\par

% structure: p0127 source=h2 target=section reason=relative-heading-rank; base=h2

\section{第六十章}

\Person{莱波里乌斯}，原为隐修士后为司铎，他依靠自身的纯洁，凭自己的自由意志和自力，而不依赖天主的助佑，开始追随白拉奇派的教理。但在受到高卢教长们告诫，并在阿非利加被\Person{奥古斯丁}纠正后，他写了一本包含其撤回声明的书，在其中他既承认自己的错误，又为所受的纠正表示感谢。同时，在纠正他对基督降生成人的错误观点时，他提出了公教的观点，承认天主圣子的独一位格，以及存在于基督内的两个性体，按其实体而存在。\par

% structure: p0129 source=h2 target=section reason=relative-heading-rank; base=h2

\section{第六十一章}

马赛的修辞学家\Person{维克托里努斯}致信其子\Person{埃泰里乌斯}，写了一部\WorkTitle{论创世纪}注释，即从卷首注释至族长亚巴郎逝世，并以诗体出版了四卷书。这些作品确有敬虔之味，但因他是一位忙于世俗文学且对圣经全无训练的人，故就思想而言，它们分量甚轻。\par

他卒于\Person{德奥多西乌斯}和\Person{瓦伦提尼亚努斯}在位时期。\par

% structure: p0132 source=h2 target=section reason=relative-heading-rank; base=h2

\section{第六十二章}

\Person{卡西安}，斯基泰族人，在君士坦丁堡由大\Person{若望}主教祝圣为执事，并在马赛成为司铎。他建立了两座隐修院，即一座男院和一座女院，至今尚存。他根据经验写作，文辞有力，或更确切地说，言中有意，语中有行。他涵盖了对各类隐修士的实用指导，著作如下：\WorkTitle{论服装}，亦\WorkTitle{论祈祷规程}，以及\WorkTitle{咏唱圣咏的习惯}（因为在埃及隐修院中，圣咏日夜咏唱），共三卷。一卷\WorkTitle{规章}，八卷\WorkTitle{论八种主要罪恶的起源、本性与补救}，每卷论一种罪恶。他还编著了与埃及教父的\WorkTitle{会谈录}，如下：\WorkTitle{论隐修士的目的与信经}、\WorkTitle{论明辨}、\WorkTitle{论三种事奉天主的圣召}、\WorkTitle{论肉身与灵性的交战}、\WorkTitle{论一切罪恶的本性}、\WorkTitle{论圣人的殉道}、\WorkTitle{论心灵的善变}、\WorkTitle{论权柄}、\WorkTitle{论祈祷的本性}、\WorkTitle{论祈祷的时长}、\WorkTitle{论成全}、\WorkTitle{论贞洁}、\WorkTitle{论天主的保护}、\WorkTitle{论灵性事物的知识}、\WorkTitle{论天主的恩宠}、\WorkTitle{论友谊}、\WorkTitle{论应当定义或不定义}、\WorkTitle{论三种古代隐修士及一种新近兴起的第四种}、\WorkTitle{论会士与隐士的目的}、\WorkTitle{论悔改的真满足}、\WorkTitle{论五旬节大斋的豁免}、\WorkTitle{论夜间幻象}、\WorkTitle{论宗徒的话：\BibleQuote{我所愿意的善，我不去行；而我所不愿意的恶，我却去作。}}、\WorkTitle{论克己}，最后，应总执事\Person{良}（后来成为罗马主教）的请求，他写了七卷书驳斥聂斯多略，\WorkTitle{论主的降生成人}。写完这部著作后，他便在马赛结束了写作与生命，时值\Person{德奥多西}和\Person{瓦伦提尼安}在位。\par

% structure: p0134 source=h2 target=section reason=relative-heading-rank; base=h2

\section{第六十三章}

\Person{菲利普}，司铎\Person{热罗尼莫}的得意门生，出版了\WorkTitle{约伯传注释}，笔法质朴。我读过他的\WorkTitle{书信集}，极其诙谐，劝勉忍受贫穷和苦难。他在\Person{马尔西安}和\Person{阿维图斯}在位期间去世。\par

% structure: p0136 source=h2 target=section reason=relative-heading-rank; base=h2

\section{第六十四章}

\Person{欧凯里}，里昂教会主教，写信给他的亲戚\Person{瓦莱里安}，\WorkTitle{论轻视世界与世俗哲学}，单篇书信，文笔显示出渊博学识与推理。他还写信给他的儿子\Person{撒洛尼}和\Person{维拉尼}（后来成为主教），一篇讨论\WorkTitle{论圣经中某些隐晦段落}。此外，他修订并缩编了圣\Person{卡西安}的某些著作，压缩为一卷，并写了其他适合教会或隐修追求的作品。他在\Person{瓦伦提尼安}和\Person{马尔西安}在位期间去世。\par

% structure: p0138 source=h2 target=section reason=relative-heading-rank; base=h2

\section{第六十五章}

\Person{文森特}，高卢人，莱林岛隐修院的司铎，勤习圣经，精于教会教义，撰写了一篇有力的辩论，文字相当精练明晰，他隐去自己的名字，题为\WorkTitle{朝圣者驳异端}。该作品第二卷大部分被盗，他重写了原作的要点，出版为一卷。他在\Person{德奥多西}和\Person{瓦伦提尼安}在位期间去世。\par

% structure: p0140 source=h2 target=section reason=relative-heading-rank; base=h2

\section{第六十六章}

\Person{西亚格里}写了\WorkTitle{论信德}，反对异端者为了破坏或取代天主圣三的名号而使用的狂妄言辞。因为他们说父不应被称为父，以免子的名号与父名谐和，而应称他为无始者、不朽者、绝对者，使得在位格上与祂有别的一切，在本性上也各自分离。这表明，本性不变的父可称为无始者，尽管圣经并未如此称祂；子的位格由祂所生，而非受造；圣神的位格由祂所发，非生亦非受造。在这位\Person{西亚格里}名下，我发现了七卷书，题为\WorkTitle{论信德与信德规则}，但由于文风不一致，我不信是由他所写。\par

% structure: p0142 source=h2 target=section reason=relative-heading-rank; base=h2

\section{第六十七章}

\Person{以撒}，安提约基雅教会司铎，其众多著作跨越很长时期，他用叙利亚文写作，尤其反对聂斯多略派和欧迪奇派。他以一首挽歌哀悼安提约基雅的沦陷，沿用了执事\Person{厄弗冷}哀悼尼科美底亚沦陷的同一格调。他在\Person{良}皇帝和\Person{马约里安}在位期间去世。\par

% structure: p0144 source=h2 target=section reason=relative-heading-rank; base=h2

\section{第六十八章}

\Person{萨尔维亚努}，马赛司铎，精通世俗与神圣文学，且毫不偏颇地说，是主教中的大师。他以经院式清晰风格撰写了许多作品，我读过以下这些：四卷\WorkTitle{论童贞的卓越}，致司铎\Person{玛策禄}；三卷\WorkTitle{论贪婪}；五卷\WorkTitle{论现世审判}；一卷\WorkTitle{论罪有应得的惩罚}，致主教\Person{撒洛尼}；另有一卷\WorkTitle{\BibleBook{训道}{篇}后段注释}，致维埃纳主教\Person{克劳狄}；一卷\WorkTitle{书信集}。他还以希腊风格创作了一卷诗，类似\OriginalTerm{六日创世}{Hexaemeron}，涵盖从创世记开始到人的创造；还有多篇向主教们发表的\WorkTitle{讲道集}，以及我确知不知多少篇\WorkTitle{论圣事}。他仍健在，年事已高。\par

% structure: p0146 source=h2 target=section reason=relative-heading-rank; base=h2

\section{第六十九章}

\Person{保利诺}撰写了\WorkTitle{论四旬期的开始}论文，我读过其中两篇：\WorkTitle{论逾越节安息日}、\WorkTitle{论服从}、\WorkTitle{论补赎}、\WorkTitle{论新教友}。\par

% structure: p0148 source=h2 target=section reason=relative-heading-rank; base=h2

\section{第七十章}

\Person{希拉里}，亚尔教会的主教，精通圣经的人，热爱贫穷，迫切渴望过简朴生活，不仅在心灵的虔诚上，也在身体的劳动上。为保持此种贫穷状态，这位出身高贵、教养迥异的人从事农耕，虽力所不及，但仍未忽略神性事务。他也是一位令人欣悦的教师，不徇情面，对众人施行训诫。他出版了一些作品，篇幅简短，却显露出不朽的才华，以及博学的头脑和雄辩的能力；其中一部对众人极有实用价值的作品，即他所著的\WorkTitle{圣奥诺拉图斯传}，为其前任。他在瓦伦蒂尼亚诺与马尔提亚诺在位期间去世。\par

% structure: p0150 source=h2 target=section reason=relative-heading-rank; base=h2

\section{第七十一章}

\Person{良}，罗马主教，致信君士坦丁堡教会的主教\Person{弗拉维亚诺}，反对司铎\Person{欧提克斯}，此人当时因觊觎主教职位，试图在教会中引入新奇事物。在此信中，他劝告弗拉维亚诺，若欧提克斯承认错误并承诺改正，就接纳他；但若他执意坚持己路，则应连同其异端一同被定罪。他在此书信中还教导并藉神圣见证证实，正如主耶稣基督应被视为圣父的真子，同样他也应被视为具有人性的真人，也就是说，他从童贞女的肉身取得了血肉身体，而非如欧提克斯所言，显示了一个来自天上的身体。他在良与马约里亚诺在位期间去世。\par

% structure: p0152 source=h2 target=section reason=relative-heading-rank; base=h2

\section{第七十二章}

\Person{莫奇穆斯}，美索不达米亚人，安提约基雅的司铎，写了一部优秀的著作\WorkTitle{驳欧提克斯}，据说正在撰写其他作品，但我尚未读到。\par

% structure: p0154 source=h2 target=section reason=relative-heading-rank; base=h2

\section{第七十三章}

\Person{提摩泰}，当\Person{普罗泰里}被亚历山大里亚人处死后，为回应民众呼声，他自愿或不自愿地允许由一个主教在彼被处死者的位置立他为主教。且为避免他因非法受任而被曾憎恨普罗泰里的民众意愿所正当罢免，他宣称邻近所有主教都是聂斯多略派，并妄图以蛮勇洗刷良心的污点，致信良皇帝，写了一部极有说服力的书，他试图用曲解的教父见证来强化此书，以致为欺骗皇帝、确立其异端，竟称罗马主教\Person{良}、加采东大公会议及所有西方主教从根本上都是聂斯多略派。但藉天主的恩宠，教会的这个敌人被驳斥并推翻在加采东大公会议上。据说他仍在流亡，依旧是异端首领，情况很可能如此。我应弟兄们之请求，为求知之故，将此书译成拉丁文，并附加了警告。\par

% structure: p0156 source=h2 target=section reason=relative-heading-rank; base=h2

\section{第七十四章}

\Person{阿斯克勒庇俄斯}，非洲人，巴盖境内一个大教区的主教，撰文反对亚略人，据说正在撰写反对多纳徒派的文章。他以即兴讲学著称。\par

% structure: p0158 source=h2 target=section reason=relative-heading-rank; base=h2

\section{第七十五章}

\Person{伯多禄}，厄德萨教会的司铎，一位著名的讲道者，写了各种主题的\WorkTitle{论集}，并仿效执事\Person{圣厄弗冷}的方式写了\WorkTitle{圣歌}。\par

% structure: p0160 source=h2 target=section reason=relative-heading-rank; base=h2

\section{第七十六章}

\Person{保禄}司铎，潘诺尼亚民族，据我亲耳听他所说，写了\WorkTitle{论守贞与轻视世俗}及\WorkTitle{生活秩序或道德矫正}，文体平实，但带有神圣的咸味。这两本书致献给一位献身于基督的贵族童贞女，名为\Person{君士坦提亚}；在其中他提及异端者\Person{若维尼安}，此人宣扬纵欲和情欲，与那种克己贞洁的生活相去甚远，以致在奢华的宴饮中结束了生命。\par

% structure: p0162 source=h2 target=section reason=relative-heading-rank; base=h2

\section{第七十七章}

\Person{帕斯托尔}主教，写了一部简短作品，以信经形式写成，其中以句子形式几乎囊括了全部教会教义。在其中，除其他无作者姓名而受他诅咒的异端外，他谴责了普里西利安派及其创始人。\par

% structure: p0164 source=h2 target=section reason=relative-heading-rank; base=h2

\section{第七十八章}

\Person{维克多}，毛里塔尼亚卡尔滕纳的主教，写了一部长篇著作反对亚略人，据我从该著作的序言所知，他通过追随者将此书送给国王\Person{盖塞里克}；还有一部著作\WorkTitle{论税吏的悔改}，在其中他依照圣经的权威为悔罪者制定了生活准则。他还写了一部安慰性的著作致一位名为\Person{巴西略}的人，题为\WorkTitle{论丧子}，充满复活的希望和良好的劝告。他还编写了许多\WorkTitle{讲道集}，这些讲道集已被整理成连续的著作，据我所知，那些关心自身救恩的弟兄们常使用它们。\par

% structure: p0166 source=h2 target=section reason=relative-heading-rank; base=h2

\section{第七十九章}

\Person{沃科尼乌斯}，毛里塔尼亚卡斯特拉努姆的主教，写了\WorkTitle{驳教会之敌：犹太人、亚略人及其他异端者}。他还撰写了一部优秀的著作\WorkTitle{论圣事}。\par

% structure: p0168 source=h2 target=section reason=relative-heading-rank; base=h2

\section{第八十章}

\Person{Musaeus}，马赛教会司铎，精通圣经，诠释精确，且文采斐然。应圣\Person{Venerius}主教之请，从圣经中选出适合一年各庆节的经文，以及适合节期的答唱咏和读经段落。教会中的读经员认为此书极有价值，既免去选段之烦忧，又有益于民众训导和礼仪庄严。他又致书给上述圣人的继任者\Person{Eustathius}主教，呈上一部宏大而卓越的著作——\WorkTitle{圣礼书}，按不同礼仪和节期分为多卷，包含读经和圣咏，供诵读和咏唱之用，且通篇充满对主的祈求和对恩惠的感谢。由此作，我们可知他是一位思想敏锐、文辞纯洁之人。据说他还曾宣讲道理，这些道理为虔诚者所珍视（据我所知），但我未曾读过。他在\Person{Leo}和\Person{Majorianus}在位时去世。\par

% structure: p0170 source=h2 target=section reason=relative-heading-rank; base=h2

\section{第81章}

\Person{Vincentius}司铎，高卢人，精通圣经，因讲道和广泛阅读而文辞优雅。他撰写了\WorkTitle{圣咏注释}。在\Place{Cannatae}，我曾听他向一位天主的人宣读此书的部分内容，同时许诺，若主保全他的生命和精力，他将以同样方式处理全部圣咏集。\par

% structure: p0172 source=h2 target=section reason=relative-heading-rank; base=h2

\section{第82章}

\Person{Cyrus}，亚历山大里亚人，职业是医生，起初是哲学家，后成为隐修士，擅长言辞。他起初优雅而有力地反对\Person{Nestorius}，但后来因过于激烈地抨击他，且依赖三段论而非圣经，开始培植\OriginalTerm{蒂莫提派}{Timothean}学说。最终他拒绝接受\Place{加采东}大公会议的决议，并认为不能默认降生成人后天主圣子包含两个性体的主张。\par

% structure: p0174 source=h2 target=section reason=relative-heading-rank; base=h2

\section{第83章}

\Person{Samuel}，\Place{厄德撒}教会司铎，据称以叙利亚文撰写了许多反对教会敌人的著作，特别是反对\Person{Nestorius}派、\Person{Eutyches}派和蒂莫提派——这些都是新兴的异端，彼此不同。为此他常提及\OriginalTerm{三兽}{triple beast}，同时以教会意见和圣经权威予以简要驳斥，向奈斯多利派指出：圣子在人性中是天主，而非单纯由童贞女所生的人；向欧提基派指出：祂具有天主所取的真实人性，而非仅是浓密空气或从天上显现的身体；向蒂莫提派指出：圣言如此成了血肉，以致圣言在本质上仍是圣言，人性仍保持人性，天主圣子的一位格由结合而非混合而成。据说他仍住在君士坦丁堡，因为在\Person{Anthemius}统治初期，我已知其著作，并知他尚在人世。\par

% structure: p0176 source=h2 target=section reason=relative-heading-rank; base=h2

\section{第84章}

\Person{Claudianus}，\Place{维也纳}教会司铎，一位杰出的演讲家和机智的辩论家。他撰写了三卷\WorkTitle{论灵魂的状况与实体}，在其中讨论除天主之外，任何事物在多大程度上是非物质的。\par

［他还写了其他一些作品，其中包括\WorkTitle{吾主受难颂歌}，开头为\Quote{Pange lingua gloriosi}。他又是维也纳主教\Person{Mamertus}的兄弟。］（见注。）\par

% structure: p0179 source=h2 target=section reason=relative-heading-rank; base=h2

\section{第85章}

\Person{Prosper}，\Place{阿基坦}人，文风博学而有力，据说撰写了多部著作。我读过一部题为\WorkTitle{编年史}的作品，记有其名，内容按圣经记述自第一人受造直到皇帝\Person{Valentinianus}去世和\Place{罗马}被汪达尔王\Person{Genseric}攻陷。我认为还有一部匿名著作（反对\Person{Cassianus}的某些作品）也出自他手；天主教会视后者为有益的，但他却斥之为有害。事实上，\Person{Cassian}和\Person{Prosper}关于天主的恩宠和自由意志的一些看法彼此相左。\Person{Leo}教宗反对\Person{Eutyches}的书信（\WorkTitle{论基督的真降生成人}）寄给多人，也被认为是由他口授而成。\par

% structure: p0181 source=h2 target=section reason=relative-heading-rank; base=h2

\section{第86章}

\Person{Faustus}，起初是\Place{莱兰}隐修院院长，后成为\Place{法国雷日}主教。他潜心研究圣经，以教会传统信经为依据，撰写了\WorkTitle{论圣神}一书，在其中藉教父的信仰表明圣神与圣父、圣子同体、同永恒，是圣三的圆满，因此是天主。他还出版了一部卓越著作\WorkTitle{论天主的恩宠——藉此我们得救}，教导天主的恩宠总是邀请、先导并扶助我们的意志，自由意志为获得虔诚效果而取得的一切并非其功绩，而是恩宠的恩赐。我还读过他的一本小册子\WorkTitle{反对亚略派和马其顿派}，在其中他主张同体的圣三；另一本则反对那些主张在受造物中有非物质事物的人，其中他以圣经见证和教父引文坚持只有天主才是非物质的。还有一封他写给小书的信，致一位名叫\Person{Graecus}的执事，此人离开大公信仰，转向了奈斯多利的邪恶。\par

在这封信中，他劝诫对方相信圣童贞玛利亚所生的并非一个后来才获得神性的单纯的人，而是真人中的真天主。他还有其他著作，但因为我未曾读过，就不想提及了。这位杰出的导师深受热情信仰和钦佩。之后，他还写信给\Person{Felix}（近卫总长，贵族身份，执政官\Person{Magnus}之子），一封非常虔诚的信，劝勉敬畏天主，这是一部极适于引导人全心悔改的作品。\par

% structure: p0184 source=h2 target=section reason=relative-heading-rank; base=h2

\section{第87章}

天主的仆人主教，撰文反对那些声称基督在世时未曾用祂的肉眼看见圣父的人——但在祂从死者中复活并升天之后，当祂被提升至天主圣父的光荣中，可以说是作为祂自我舍弃的赏报和祂殉道的补偿时。在这部著作中，他既凭自己的论证，也凭圣经的见证，指出主耶稣从祂因圣神受孕和童贞玛利亚诞生——由此真天主在真人内，同时祂自己也是人而成为天主——开始，便因天主与人的特殊而完全的合一，常用祂的肉眼看见圣父和圣神。\par

% structure: p0186 source=h2 target=section reason=relative-heading-rank; base=h2

\section{第八十八章}

\Person{维多略}，阿基坦人，一位精于计算者，应罗马主教\Person{圣希拉略}之请，以最审慎的研究追随其四位前人（即\Person{希玻律托}、\Person{欧瑟伯}、\Person{德奥斐洛}和\Person{普罗斯柏}），编撰了一个\WorkTitle{逾越节周期}，并将年数延续至五百三十二年，如此计算：在533年，逾越节庆典应再次发生于与主受难和复活发生的那第一年相同的月份、日子和月相。\par

% structure: p0188 source=h2 target=section reason=relative-heading-rank; base=h2

\section{第八十九章}

\Person{德奥多勒}，\Place{居鲁士}主教（因波斯王\Person{居鲁士}所建之城至今在叙利亚仍保留其创立者之名），据说撰写了多部著作。我所知的有以下：\WorkTitle{论主的降生成人}、\WorkTitle{驳司铎欧提克斯和\Place{亚历山大里亚}主教\Person{狄奥斯哥鲁}}——他们否认基督具有人性肉身；这些强有力的著作藉理性和圣经的见证证明，祂从童贞母亲所获得的母性实质中拥有真实的肉身，正如祂拥有从天主圣父永恒生育所接受的真神性一样。另有十卷本\WorkTitle{教会史}，他模仿\Place{凯撒勒雅}的\Person{欧瑟伯}而作，始于\Person{欧瑟伯}结束之处，延至他本人的时代，即从\Person{君士坦丁}的二十年统治至\Person{老良}登基，他在\Person{良}在位期间去世。\par

% structure: p0190 source=h2 target=section reason=relative-heading-rank; base=h2

\section{第九十章}

\Person{甄纳迪}，\Place{君士坦丁堡}教会宗主教，口才出众、才智卓越，藉对古人的阅读而装备丰富，得以逐字注释全部\BibleBook{达尼尔}{先知书}。\par

他也撰写了许多讲道词。他在\Person{老良}皇帝在位时去世。\par

% structure: p0193 source=h2 target=section reason=relative-heading-rank; base=h2

\section{第九十一章}

\Person{德奥杜禄}，\Place{科埃莱叙利亚}的一位司铎，据说撰写了多部著作，但传到我这的只有一部，是他所著的\WorkTitle{论圣经的和谐}，即新旧约圣经，驳斥古代异端者——他们因礼仪规诫中的差异而声称旧约的天主不同于新约的天主。在这部著作中，他表明：由于同一天主——两部圣经的作者——的经世安排，古老的律法藉梅瑟以祭献礼仪和司法性法律赐给古人，而另一部律法则藉基督的来临在神圣奥迹和未来许诺中赐给我们；它们不应被视为不同，而是由同一圣神和同一作者所吩咐，因为这些事若仅按文字遵守就会杀害，若按精神遵守则会使心灵得生命。这位作者在三年前\Person{芝诺}皇帝在位时去世。\par

% structure: p0195 source=h2 target=section reason=relative-heading-rank; base=h2

\section{第九十二章}

[\Person{西多尼}，\Place{阿维尔尼}人主教，撰写了几部可接受的作品，他本人教义纯正，博学于神学与人文，才华卓越，写作了一卷相当可观的\WorkTitle{书信}，致不同人士，以各种诗律或散文写成，显示了他的文学才能。在野蛮残暴肆虐\Place{高卢}的时期，他仍以基督徒的强韧著称，被视为公教教父和杰出圣师。他在\Person{良}与\Person{芝诺}统治的风暴时期活跃。]\par

% structure: p0197 source=h2 target=section reason=relative-heading-rank; base=h2

\section{第九十三章}

\Person{安提约基雅的若望}，先为文法学家，后为司铎，撰文反对那些声称基督仅在单一实体中受钦崇、不承认在基督内应辨认两个性体的人。他按照圣经记载教导：在祂内，天主和人在一个位格中存在，而非肉体和圣言在一个性体中。\par

他也攻击了\Place{亚历山大里亚}主教\Person{济利禄}的某些意见，这些意见是\Person{济利禄}不智地针对\Person{聂斯多略}而发表的，如今却成为提摩替派信徒的鼓励和力量。据说他仍在世并宣讲。\par

% structure: p0200 source=h2 target=section reason=relative-heading-rank; base=h2

\section{第九十四章}

[\Person{哲拉修}，罗马主教，撰写了\WorkTitle{驳欧提克斯和聂斯多略}——一部伟大而著名的著作，以及论述圣经和圣事各部分的\WorkTitle{论著}，风格典雅。他也撰写了\WorkTitle{驳伯多禄和阿卡西乌书函}，这些至今保留在天主教会内。他还仿照主教\Person{盎博罗削}的方式撰写了\WorkTitle{赞美诗}。他在\Person{阿纳斯塔修斯}皇帝在位时去世。]\par

% structure: p0202 source=h2 target=section reason=relative-heading-rank; base=h2

\section{第九十五章}

\Person{何诺拉特}，非洲\Place{君士坦提纳}主教，致信给一位名叫\Person{阿卡迪乌}的人——此人因宣认公教信仰而被\Person{耿塞里克}王放逐至非洲。这封信劝勉他为基督忍受艰苦，并以近代实例和圣经例证坚固他，表明在信仰宣认上的坚忍不仅涤除以往之罪，也获得殉道的祝福。\par

% structure: p0204 source=h2 target=section reason=relative-heading-rank; base=h2

\section{第九十六章}

\Person{塞雷亚利斯}主教，非洲裔，应亚略派主教\Person{马克西慕}之请，询问他是否可以用少数几项圣经典证、无需任何争议性断言来建立公教信仰。他奉主名而行，真理自身帮助他，不仅用少数证言如马克西慕嘲弄般所求，而且用新旧约大量的证明经文来证实，并将之出版为一本小册子。\par

% structure: p0206 source=h2 target=section reason=relative-heading-rank; base=h2

\section{第九十七章}

\Person{欧任纽}，非洲迦太基主教、公开的明认者，受汪达尔王\Person{亨塞里克}之命撰写一份天主教信仰的阐明，并特别讨论\OriginalTerm{同一本体}{Homoousian}一词的含义。他征得仍持守天主教信仰的毛里塔尼亚、非洲、撒丁岛及科西嘉岛所有主教与明认者的同意，撰写了一部信仰之书，不仅有圣经引文作支持，更有教父们的见证，并遣其同难伴侣呈递。如今，因他忠诚的口舌而被放逐，他如焦虑的牧人照看羊群一般，留下了著作，敦促他们务必铭记信仰与唯一神圣的圣洗圣事，甚至不惜一切代价。他还写下了他与亚略派领袖们通过信使进行的\WorkTitle{讨论}，并将其送交其管家转呈\Person{亨塞里克}。同样，他也向同一人呈递了为基督徒和平请愿的\WorkTitle{辩护}性质的文件，据称他至今仍在世，为教会的巩固而存。\par

% structure: p0208 source=h2 target=section reason=relative-heading-rank; base=h2

\section{第九十八章}

\Person{波默里乌斯}，原是毛里塔尼亚人，在高卢被祝圣为司铎。他撰写了一部八卷的辩证论著\WorkTitle{论灵魂及其属性}，以及一部\WorkTitle{论复活}及其对今世信徒的特殊意义、对所有人的普遍意义的著作，文笔清晰、风格优雅，采用\Person{朱利安}主教与\Person{维鲁斯}司铎之间的对话形式。第一卷讨论灵魂是什么，以及它在何种意义上被视为按天主的肖像受造；第二卷讨论灵魂应被视作有形还是无形；第三卷讨论第一人的灵魂如何受造；第四卷讨论出生时进入身体的灵魂是新造而无罪的，抑或如从根部长出的嫩芽般从第一人的实体产生，因而也带来了第一人的原罪；第五卷回顧前四卷的讨论，探讨灵魂的能力（即其潜能），并指出其潜能来自单一而纯粹的意志；第六卷讨论宗徒所说的肉性与灵性的冲突从何而来；第七卷讨论肉性与灵性在生命、死亡和复活方面的区别；第八卷回答关于世界终结时被预言将要发生之事的问题，即通常就复活所提出的问题。我记得曾读过他写给一位名叫\Person{普林奇皮乌斯}的人的劝勉之作\WorkTitle{论轻视世界与可朽之物}，以及另一部题为\WorkTitle{论恶习与美德}的著作。据说他还写了其他作品，只是我未曾得见，并且他仍在写作。他至今在世，其生活堪配基督徒身份及其在教会中的品级。\par

% structure: p0210 source=h2 target=section reason=relative-heading-rank; base=h2

\section{第九十九章}

[我，\Person{根纳狄}，马赛的司铎，撰写了八卷\WorkTitle{驳一切异端}、五卷\WorkTitle{驳聂斯脱利}、十卷\WorkTitle{驳欧提基}、三卷\WorkTitle{驳伯拉纠}，以及论著\WorkTitle{论千年王国}和\WorkTitle{论圣若望默示录}，还有一封致真福\Person{格拉修}（罗马主教）的书信\WorkTitle{论我的信经}。]\par

\SourceRecord{Translated by Ernest Cushing Richardson.；Nicene and Post-Nicene Fathers, Second Series；Vol. 3.；Edited by Philip Schaff and Henry Wace.；Buffalo, NY: Christian Literature Publishing Co.,；1892.；<http://www.newadvent.org/fathers/2719.htm>.}
